\scanpage{135}
% lines 1–3
\deleted{d) Considérons alors la factorisation de Stein.
\[
f\qquad X\longrightarrow Y'\unsure{-|\to}Y
\]}
\unsure{Cela|Ceci} nous ramène au cas où $Y$ est \uline{réduit}.

% lines 4–12
c) \unsure{Prenons} la factorisation de Stein
\[
X-Y'-Y\deleted{\unsure{'}}
\]
où $Y'-Y\deleted{\unsure{'}}$ est fini surjectif, donc un épimorphisme fini factorisable en produit d'épimorphismes stricts. Par ``Seshadri'', $\underline{\mathrm{Pic}}_{Y}-\underline{\mathrm{Pic}}_{Y'}$ est affine. On est ramené \deleted{à prouver que $\underline{\mathrm{Pic}}_{Y'}-\underline{\mathrm{Pic}}_X$ l'est, \unsure{ou}} ce qui nous ramène au cas $f_{*}(\underline O_X)\simeq \underline O_Y$

% lines 13–15
d) Mais dans ce cas on sait aussitôt que $\underline{\mathrm{Pic}}_X-\underline{\mathrm{Pic}}_Y$ est un \unsure{mono}, donc \deleted{\unsure{un noyau est}} une immersion fermée, donc affine

% slanted
\margin{Attention au fait que $\underline{\mathrm{Pic}}_X$ et $\underline{\mathrm{Pic}}_Y$ ne sont pas de type fini sur $k$, \unsure{on aura d'inf|on aura d'of} donc pas \unsure{néc$^{t}$} affine [ex. $\mathbb Z\to G$, $G=\unsure{\mathrm G_m|G_a|\mathrm{Grass}}$ ou \unsure{courbes|$G_a$} ou VA]}

% lines 16–18
NB. On s'est servi de l'existence des $\underline{\mathrm{Pic}}_Z$ pour $Z=Y_{\mathrm{réd}}$, $X_{\mathrm{réd}}$, et pour les préschémas finis sur $Y_{\mathrm{réd}}$ ; \deleted{\unsure{Est}} on peut sans doute s'en passer \unsure{\textemdash]|(?)\textnormal{]}}.
% archive sheet number: 134

